% Clase 19 --- El juego epsilon/delta y el entorno reducido (§1.3).
% El docente: "si un usuario me da un epsilon esto me define una banda
% horizontal, una banda de precisión... el juego es hallar un delta bastante
% pequeño, pero no nulo, tal que la imagen de la banda de radio delta esté
% adentro de la banda horizontal... yo no podría tomar delta acá, porque esta
% parte de la función se va afuera de la banda". Y además: "nunca se examina
% el valor de f en el punto x sub cero".
\begin{center}
\begin{tikzpicture}
\begin{axis}[
    calcfig,
    axis lines=left,
    width=12cm, height=7cm,
    domain=1.2:6.8, samples=100,
    xmin=0.8, xmax=7.6, ymin=1.4, ymax=5.4,
    xtick={4}, xticklabels={$x_0$},
    ytick={3}, yticklabels={$L$},
    xlabel={$x$}, ylabel={$y$},
    xlabel style={font=\small, at={(axis description cs:1,0)}, anchor=north},
    ylabel style={font=\small, at={(axis description cs:0,1)}, anchor=east},
  ]
  % banda horizontal de precisión: L +- epsilon
  \fill[figamber!18] (axis cs:0.8,2.4) rectangle (axis cs:7.05,3.6);
  % banda vertical que SÍ funciona: radio delta
  \fill[figblue!14] (axis cs:3.1,1.4) rectangle (axis cs:4.9,5.05);
  % radio demasiado grande: la imagen se escapa de la banda de precisión
  \draw[figred, dashed, line width=0.8pt] (axis cs:1.8,1.4) -- (axis cs:1.8,5.05);
  \draw[figred, dashed, line width=0.8pt] (axis cs:6.2,1.4) -- (axis cs:6.2,5.05);

  \addplot[figblue, line width=1.1pt] {3 + 0.6*(x-4) + 0.08*(x-4)^2};

  % guías del punto y del valor irrelevante f(x_0)
  \draw[figaux] (axis cs:0.8,3) -- (axis cs:4,3);
  \node[figptoabierto] at (axis cs:4,3) {};
  \node[figpto, fill=figred, minimum size=5.5pt] at (axis cs:4,4.6) {};
  \node[figlbl, figred, right=3pt] at (axis cs:4,4.6) {$f(x_0)$};

  % rótulos
  \node[figlbl, figamber, right] at (axis cs:7.08,3.42) {$L+\varepsilon$};
  \node[figlbl, figamber, right] at (axis cs:7.08,2.58) {$L-\varepsilon$};
  \node[figlblaux, align=center] at (axis cs:4,1.75)
    {radio $\delta_\varepsilon$\\ (sirve)};
  \node[figlbl, figred, align=center, anchor=south] at (axis cs:6.2,5.10)
    {radio demasiado grande};
\end{axis}
\end{tikzpicture}\\[0.5em]
{\small\itshape Una precisión $\varepsilon$ define la banda horizontal
$L\pm\varepsilon$; el juego consiste en devolver un radio $\delta_\varepsilon$
tal que la imagen de la banda vertical caiga entera adentro de la horizontal.
El radio punteado en rojo fracasa: parte de la gráfica se escapa. El valor
$f(x_0)$ puede estar en cualquier lado --- el entorno es \emph{reducido} y
nunca se lo mira.}
\end{center}
